% ============================================================================
% Compiles with the standard `article` class (no special packages needed), so it
% builds anywhere (Overleaf, texlive, etc.). To submit to an ACM venue, replace
% this preamble block with the acmart version noted in README.md.
% ============================================================================
\documentclass[10pt,twocolumn]{article}

\usepackage[margin=0.75in]{geometry}
\usepackage{graphicx}
\usepackage{booktabs}
\usepackage{xcolor}
\usepackage{amsmath}
\usepackage[hidelinks]{hyperref}
\usepackage{authblk}
\usepackage{tikz}
\usetikzlibrary{arrows.meta,positioning,shapes.geometric,backgrounds,fit}

\renewcommand\Authfont{\normalsize}
\renewcommand\Affilfont{\small\itshape}
\setlength{\affilsep}{0.3em}

\title{\textbf{Connection-Churn Interference in Shared-Dataplane Service Meshes:
Measurement and Mitigation}}

\author[]{Ajitesh Kumar Singh}
\author[]{Shashank Tippanavar}
\affil[]{IIIT Bangalore}
\date{}

\begin{document}
\twocolumn[{%
\maketitle
\begin{center}
\begin{minipage}{0.86\textwidth}
{\centering\textbf{Abstract}\par}\medskip
\small
Istio's ambient mode reduces the cost of a service mesh by replacing the per-pod sidecar
proxy with a single shared proxy (\emph{ztunnel}) per node. This design improves resource
efficiency but introduces a new failure mode: because all pods on a node share one proxy,
a single workload can degrade co-located neighbors through the proxy alone, without
saturating their CPU, their application code, or the network.

We study this shared-proxy interference on a standard microservice benchmark
(DeathStarBench Social Network) running on a managed Kubernetes cluster with Cilium and
Istio ambient, and report three findings. First, two distinct causes---connection
\emph{churn} (many short-lived, TLS-bearing connections) and shared application-backend
saturation---produce the same victim symptom (median latency remains stable while the
99th-percentile tail inflates by roughly $17\times$) yet leave opposite signatures inside
the proxy. Second, proxy saturation behaves as a threshold: it engages sharply and
provides no early warning, so an effective mitigation must react rather than predict.
Third, the scarce resource is the proxy's connection-handling capacity, which conventional
QoS mechanisms do not observe. We show that Cilium's Bandwidth Manager, the closest
existing mechanism, cannot mitigate connection churn even under an aggressive cap, because
it paces bytes while a large number of small connections remain within any reasonable byte
budget.

Guided by these findings, we build a closed-loop controller that monitors a single proxy
saturation signal and, when the proxy is overloaded, paces the aggressor's new-connection
rate. Over $n{=}12$ trials it reduces victim tail latency from $47\,$ms to $3.7\,$ms
($12.7\times$, within 30\% of the interference-free baseline; Mann--Whitney
$p{<}10^{-4}$) while continuing to serve the aggressor, adds $9\,$ns per packet at the
enforcement primitive, and imposes no measurable overhead when the proxy is unsaturated. The underlying principle generalizes: in a
shared-dataplane mesh, measure the resource that is actually scarce, at the point where it
is scarce, and control the input that consumes it.
\end{minipage}
\end{center}
\vspace{1.2em}
}]

\section{Introduction}
A service mesh provides a dedicated infrastructure layer for microservice
communication---encryption, retries, and telemetry---so that these concerns are handled
outside application code. The conventional design places a proxy \emph{sidecar} alongside
every pod. This gives each pod an isolated proxy but is resource-intensive: a large
deployment runs hundreds of proxy instances, each consuming memory and CPU.

Istio's \emph{ambient} mode removes the sidecars. Each node instead runs a single shared
proxy, \emph{ztunnel}, that handles the layer-4 traffic (mutual TLS, telemetry) for all
pods on that node. This improves resource efficiency and is a primary reason ambient mode
is being adopted. However, sharing a proxy also shares a failure domain: if one pod
overloads its node's ztunnel, every other pod routed through it can be affected---including
pods that are otherwise idle and well within their own resource limits.

This paper studies that effect and asks three questions, each answered with measurements
on a running system:
\begin{enumerate}
\item \emph{Does it occur?} Can one workload degrade a co-located neighbor purely through
the shared proxy, with the neighbor's own resources unsaturated? We show that it can, and
we separate this from a similar-looking cause (shared application backends).
\item \emph{What signal exposes it, and does it provide early warning?} We find that the
proxy's own telemetry detects and classifies the condition almost perfectly, but only as
it occurs, not in advance.
\item \emph{Can it be mitigated, and why do existing tools not already do so?} A small
reactive controller mitigates it; existing tools miss it because they observe and control
the wrong resource.
\end{enumerate}

Our central observation is a mismatch between what conventional infrastructure tools
control and what is actually scarce. Autoscalers add application CPU; bandwidth managers
pace bytes. But the resource a shared mesh proxy exhausts first, under an adversarial or
misbehaving neighbor, is its capacity to \emph{establish connections}: each new mutual-TLS
handshake is comparatively expensive in CPU while carrying almost no application bytes. A
workload that opens many short-lived connections (\emph{connection churn}) can therefore
saturate the shared proxy while appearing benign to any byte- or application-CPU-based
tool. We demonstrate this directly---Cilium's Bandwidth Manager, which paces bytes, does
not mitigate churn---and show that pacing the correct quantity, the new-connection rate,
triggered by a proxy-saturation signal, restores victim tail latency by $12.7\times$.

\paragraph{Contributions.}
\begin{itemize}
\item A controlled characterization, on a standard microservice benchmark, showing that
shared-proxy interference is real and separating it from application-backend
contention---two causes with an identical victim symptom but opposite proxy signatures
(\S\ref{sec:char}).
\item An empirical signal study showing the proxy's telemetry is a near-perfect
\emph{binary detector} and \emph{mechanism classifier}, but not a graded predictor and
not an early-warning signal---which dictates a reactive, not predictive, controller
(\S\ref{sec:signal}).
\item A closed-loop connection-rate controller that restores victim tail latency
$12.7\times$, with a $9\,$ns/packet data-plane cost and a true no-op when healthy
(\S\ref{sec:mitigation}).
\item A head-to-head result showing why the closest existing mechanism (byte-rate
pacing) cannot address this interference class, and why the usual escape hatch
(autoscaling) structurally does not apply to a shared proxy (\S\ref{sec:diff}).
\end{itemize}

\section{Background and Setup}
\label{sec:setup}
\paragraph{Ambient mode.} In Istio ambient mode, a per-node DaemonSet
process (ztunnel, written in Rust on the Tokio async runtime) transparently intercepts
the L4 traffic of every meshed pod on its node and tunnels it over mutual TLS. There is
exactly one ztunnel per node and it is not something an application can scale
out---unlike an optional per-namespace ``waypoint'' proxy, a tenant cannot spin up more
ztunnels to escape contention. That singleton, always-in-the-path property is exactly
what makes it worth studying.

\paragraph{Testbed.} We run on Amazon EKS (Kubernetes 1.33) with non-burstable
\texttt{c6i.2xlarge} nodes (8 vCPU), Cilium 1.20 as the CNI (overlay mode), and Istio
1.31 in ambient mode. The workload is the DeathStarBench Social Network benchmark (27
microservices: C++/Thrift services with MongoDB, Redis, and memcached backends). We
drive it with the open-loop \texttt{wrk2} generator (which avoids coordinated omission)
and report latency percentiles from full HDR histograms. Our victim is the
\texttt{home-timeline} read endpoint held at a fixed 100 requests/s.

\paragraph{Isolating the proxy.} To attribute effects to the shared proxy rather than to
a full node, we (i) co-locate all workloads and the ztunnel on one ``under-test'' node,
(ii) \emph{cap ztunnel's CPU at one core} so it saturates while the 8-core node still has
plenty of headroom, and (iii) run the load generator on a separate node. At every data
point we report node CPU and the ztunnel's own CPU/throttling so a reader can see that
the node is not the bottleneck. Because ztunnel is a distroless image, we read its cgroup
CPU counters from the node and scrape its Tokio runtime metrics over the network.

\paragraph{A note on rigor.} Latency is heavy-tailed and runs vary, so every headline
number is a mean over repeated trials with the spread reported; load-level order is
randomized across trials; and we lead with the strongest, most reproducible framing of
each result. Where a result is a smaller pilot ($n{=}3$) rather than a full sweep, we say
so.

\paragraph{On the eBPF enforcer and Cilium.}
\label{sec:setup-ebpf}
Our data-path enforcer is an eBPF program attached to the aggressor pod's virtual
interface that rate-limits connection-initiating (SYN) packets via a token bucket in a
BPF map that the control loop updates. It compiles and load-verifies, and we benchmark it
in isolation (\S\ref{sec:diff}). One integration wrinkle is worth stating up front so our
numbers are not over-read: Cilium attaches its own program to the same interface via the
newer \emph{tcx} hook, which executes ahead of the classic TC hook, so a program attached
the classic way never sees the packets. Attaching ours with explicit tcx ordering (ahead
of Cilium's) is the correct production path but requires tooling we treat as engineering.
To keep the \emph{measured} results honest we therefore enforce the identical
new-connection-rate policy at the connection-admission layer for all closed-loop
experiments; since both enforce the same cap, they are equivalent for the victim-facing
outcome, and we flag the eBPF-in-CNI integration as the one piece not yet wired
end-to-end.

\begin{table}[t]
\centering
\small
\caption{Victim 99th-percentile latency across conditions. The shared-proxy attack
inflates the tail $\sim$16$\times$; byte-rate pacing does not mitigate it; connection-rate
pacing restores it. Values are mean with 95\% bootstrap CI. The core comparison
(baseline/unmitigated/ours) uses $n{=}12$ trials; the two Bandwidth Manager rows use
$n{=}3$.}
\label{tab:summary}
\begin{tabular}{lcr}
\toprule
Condition & $n$ & Victim P99 (ms) \\
\midrule
No aggressor (baseline)              & 12 & $2.9$ \,[2.8, 2.9] \\
Connection-churn attack, unmitigated & 12 & $47.0$ [36, 57] \\
\;+ Cilium Bandwidth Manager (1\,Mbit)   & 3 & $56.5 \pm 3.2$ \\
\;+ Cilium Bandwidth Manager (256\,kbit) & 3 & $29.2 \pm 13.1$ \\
\;+ \textbf{Ours (connection-rate pacing)} & 12 & $\mathbf{3.7}$ \,[3.6, 3.8] \\
\bottomrule
\end{tabular}
\end{table}

\section{Evaluation Overview}
\label{sec:eval}
Our evaluation answers five questions, each in a following section:
\textbf{(Q1)} Does interference occur purely through the shared proxy, with the victim's own
resources unsaturated, and can it be distinguished from application-backend contention?
(\S\ref{sec:char})
\textbf{(Q2)} Which proxy signal exposes the condition, and does it provide early warning?
(\S\ref{sec:signal})
\textbf{(Q3)} Does a reactive connection-rate controller restore victim tail latency while
continuing to serve the aggressor? (\S\ref{sec:mitigation})
\textbf{(Q4)} Do existing mechanisms (byte-rate pacing, autoscaling) address this
interference? (\S\ref{sec:diff})
\textbf{(Q5)} What is the overhead, and how sensitive is the controller to its parameters?
(\S\ref{sec:overhead})
All experiments use the testbed and methodology of \S\ref{sec:setup} (managed EKS, Cilium
1.20, Istio 1.31 ambient; victim held at 100\,req/s; open-loop \texttt{wrk2}; HDR
histograms; repeated trials with randomized ordering).

\section{Shared-Proxy Interference and Its Signatures}
\label{sec:char}
We inject two kinds of aggressor, both co-located on the victim's node, and sweep their
intensity. The first is a \emph{connection-churn} aggressor that repeatedly opens and
closes connections to a dummy service that shares \emph{only} the ztunnel with the victim
(no shared application backend). The second is a \emph{compose-post} aggressor: a real
DeathStarBench write workload that shares application backends (post-storage, user,
Redis) with the victim's read path.

\begin{figure}[t]
\centering
\includegraphics[width=\linewidth]{figs/fig_two_mechanisms.pdf}
\caption{Two aggressors produce the \emph{same} victim symptom (a) but leave
\emph{opposite} fingerprints in the proxy (b). Connection churn saturates ztunnel's
capped CPU; the backend aggressor degrades the victim while leaving ztunnel nearly idle.
Error bars are $\pm$1 SD over 5 (churn) / 3 (compose) trials.}
\label{fig:two}
\end{figure}

\paragraph{The victim symptom.} Figure~\ref{fig:two}(a) shows the effect. With no
aggressor, the victim's P99 sits at $2.9\,$ms. Under connection churn it jumps to
$\sim$50\,ms---a $17$--$18\times$ tail amplification---while the victim's \emph{median}
latency barely moves (it stays near $2\,$ms). This ``flat median, exploding tail''
signature is exactly why the problem is easy to miss: a dashboard watching average
latency sees nothing.

\paragraph{It is the proxy, not the node.} Throughout the churn sweep the under-test node
sits at only $\sim$30\% CPU, while ztunnel is pinned at its 1-core cap and is actively
CPU-throttled ($\sim$680 throttle events per minute). The victim's own service is
untouched. The only shared, saturated resource is the proxy---so the interference is
attributable to it and not to node exhaustion or application code.

\paragraph{A second, distinct cause.} The compose-post aggressor produces an
almost identical victim symptom (Figure~\ref{fig:two}a, blue): no effect until a
threshold, then a violent tail blow-up ($95\times$ at high load). But its proxy
fingerprint is the \emph{opposite} (Figure~\ref{fig:two}b): ztunnel stays near-idle
($<0.1$ core, zero throttling). Here the bottleneck is a shared application backend
saturating, not the proxy. This matters for two reasons. It is a confound we must rule
out to make an honest ``shared-proxy'' claim---which the disjoint-backend churn aggressor
does. And it foreshadows the signal question: the victim cannot tell the two apart, but,
as we show next, the proxy can (Figure~\ref{fig:sig}).

\begin{figure}[t]
\centering
\includegraphics[width=0.82\linewidth]{figs/fig_signature.pdf}
\caption{The same victim tail arises from two regimes with opposite proxy-CPU
signatures. Proxy CPU is diagnostic of \emph{which} problem is occurring, not of
\emph{how bad} the victim tail is.}
\label{fig:sig}
\end{figure}

\paragraph{It is a threshold, not a dial.} A practically important detail:
the churn effect is a \emph{step}. Even our lowest churn setting already pins ztunnel at
its cap, so pushing harder does not make things proportionally worse---the proxy is
simply saturated or it is not. This shapes everything that follows: if the problem has no
gentle ramp, there is little to predict, and a good controller should react quickly at
the threshold rather than try to forecast severity.

\section{Signal Selection}
\label{sec:signal}
A mitigation needs something to watch. We recorded ztunnel's observable
signals---connection rate, worker busy-time, in-flight task count, CPU throttling---at
high cadence (every $\sim$200\,ms) across five trials while an aggressor switched on
mid-trace, and asked three questions of each signal.

\paragraph{(1) Does it track how \emph{bad} the tail is?} No. Within the saturated
regime (excluding the on/off transition), the correlation between any proxy signal and
the victim's P99 is essentially zero (worker busy-rate: $r={-}0.06\pm0.27$; in-flight
tasks: $r={-}0.38\pm0.15$). This is the honest consequence of the threshold in
\S\ref{sec:char}: once the proxy is saturated, its telemetry does not encode \emph{how
much} the victim is hurting. We therefore do \emph{not} claim a predictive signal.

\paragraph{(2) Does it detect \emph{whether} the proxy is saturated?} Almost perfectly.
Figure~\ref{fig:detector} shows the worker busy-rate is cleanly bimodal: idle
(median $0.0$, 95th pct $0.05$) versus loaded (median $2.0$, 5th pct $1.8$)---a $38\times$
gap with no overlap. A single fixed threshold separates the two states with 100\%
accuracy across $\sim$1000 samples. So the proxy is a superb \emph{binary detector} of
its own saturation.

\paragraph{(3) Can it tell the two mechanisms apart?} Yes. CPU throttling is nonzero only
under proxy (churn) saturation and stays exactly zero under backend saturation, while
in-flight task count rises in both. So a simple rule---\emph{throttling $>0$ means the
proxy is the problem; throttling $=0$ but tasks rising means a backend is}---classifies
the mechanism from the proxy's telemetry alone (Figure~\ref{fig:classify}).

\begin{figure}[t]
\centering
\resizebox{0.92\linewidth}{!}{\input{figs/body_classifier}}
\caption{Classifying the cause from proxy telemetry. The victim symptom is identical for
both mechanisms, but two proxy signals separate them, and each points to a different
remedy.}
\label{fig:classify}
\end{figure}

\paragraph{(4) Does it warn us early?} No. Comparing signal onset to victim-tail onset
across trials, the signal moves essentially \emph{simultaneously} with the tail (within
$\sim$1\,s, and if anything slightly after). Combined with the threshold shape, this
rules out predictive control and points to a fast \emph{reactive} loop with hysteresis.

\begin{figure}[t]
\centering
\includegraphics[width=0.72\linewidth]{figs/fig_detector.pdf}
\caption{The proxy's worker busy-rate is a clean on/off detector of saturation: idle and
loaded distributions do not overlap, so a single threshold separates them perfectly.}
\label{fig:detector}
\end{figure}

The practical upshot is a controller design, not a predictor: \emph{detect saturation
fast via a threshold on proxy busy-rate/throttling; react; use hysteresis to avoid
flapping}.

\section{A Reactive Connection-Rate Controller}
\label{sec:mitigation}

\begin{figure}[t]
\centering
\resizebox{\linewidth}{!}{\input{figs/body_controlloop}}
\caption{System overview. The aggressor and victim share one node-level proxy. A
userspace controller monitors the proxy's saturation signal (dashed) and, when the proxy
is overloaded, caps the aggressor's new-connection rate (red); the victim and all other
traffic are untouched. The proxy cannot be scaled out (one per node), so pacing the input
is the available lever.}
\label{fig:arch}
\end{figure}

Figure~\ref{fig:arch} shows the system. The remainder of this section explains the choice
of control input (\S\ref{sec:mitigation}), the controller itself, and the result.

\paragraph{Why connections, not bytes.} \S\ref{sec:char}--\ref{sec:signal} tell us the
scarce resource is the proxy's capacity to \emph{establish} connections (each mutual-TLS
handshake costs CPU; the bytes are negligible). So the input to control is the
aggressor's \emph{new-connection rate}. We first confirm this is the right lever
(Figure~\ref{fig:principle}): as we cap the aggressor's new-connection rate, there is a
clean knee---once the cap drops below what the proxy can absorb, ztunnel stops throttling
and the victim's tail collapses back to baseline. Capping connections un-saturates the
proxy; nothing about bytes is involved.

\begin{figure}[t]
\centering
\includegraphics[width=0.82\linewidth]{figs/fig_principle.pdf}
\caption{Capping the aggressor's new-connection rate un-saturates the proxy. Below the
knee ($\sim$800 conns/s here) ztunnel stops throttling (right axis) and the victim tail
returns to baseline (left axis).}
\label{fig:principle}
\end{figure}

\paragraph{The controller.} We build a small closed loop with two interchangeable
enforcement backends. A userspace agent polls the proxy's saturation signal ($\sim$1\,s);
when busy-rate crosses a high threshold it caps the aggressor's new-connection rate to a
setpoint; it releases only after the proxy has stayed healthy for a dwell period. The
\emph{policy} (detect $\to$ cap $\to$ dwell) is identical across backends; only the
\emph{enforcement point} differs. We implement two: (i) a connection-layer token bucket
on new-connection admission, which is what produces the closed-loop numbers in this
section and \S\ref{sec:diff}; and (ii) an eBPF program that rate-limits connection-
initiating (SYN) packets in the TC data path, which is the production-intended enforcer.
We are explicit about which backend produced which number: the closed-loop
victim-recovery results use backend (i); the $9\,$ns/packet cost in \S\ref{sec:diff} is a
microbenchmark of backend (ii) in isolation. Because both enforce the same
new-connection-rate cap, they are functionally equivalent for the victim-facing result;
the remaining gap---attaching (ii) \emph{ahead of} Cilium's own data-path hook---is an
integration detail we discuss in \S\ref{sec:setup-ebpf} and scope as engineering, not a
difference in mechanism.

\paragraph{A control subtlety worth naming.} Our first controller \emph{oscillated}. The
reason is instructive: pacing \emph{works}, so as soon as it un-saturates the proxy the
naive ``recovered $\rightarrow$ release'' rule immediately lifts the cap, the proxy
re-saturates, and the cycle repeats every few seconds. The action erases the very signal
that triggered it. The fix is standard once you see it---a dwell time plus a
sustained-recovery requirement---after which the controller settles into a single clean
pacing episode. This is a concrete instance of the general lesson that a controller acting
on a self-suppressing signal must be hysteretic.

\begin{figure}[t]
\centering
\includegraphics[width=0.62\linewidth]{figs/fig_mitigation.pdf}
\caption{Closed-loop result ($n{=}12$ per bar, 95\% bootstrap CI). The controller reduces
the victim's P99 from $47\,$ms to $3.7\,$ms---within 30\% of the interference-free
baseline---while still serving the aggressor. The wide interval on the unmitigated bar
reflects the intrinsic variance of proxy saturation, not measurement error.}
\label{fig:mit}
\end{figure}

\paragraph{Result.} Figure~\ref{fig:mit} shows the outcome, over $n{=}12$ trials per
condition with 95\% bootstrap confidence intervals. Unmitigated, the victim's P99 is
$47.0\,$ms (95\% CI $[36.1, 56.6]$). With the controller enabled it is $3.72\,$ms (95\% CI
$[3.64, 3.79]$)---a $12.7\times$ reduction that recovers 98\% of the interference gap and
lands within 30\% of the no-aggressor baseline ($2.88\,$ms), while the aggressor still
receives service at the setpoint (it is paced, not shut off). A Mann--Whitney $U$ test
confirms the improvement is significant ($p{=}2.8\times10^{-5}$; we use a non-parametric
test because the unmitigated distribution is not normal). The mitigated condition is
tightly reproducible (CV $4\%$); the unmitigated condition is intrinsically noisy
(CV $40\%$) because the proxy is pinned at its saturation cap and scheduling jitter
dominates---a property of the phenomenon, not of measurement, and the reason the tail
rather than the mean is the meaningful metric.

\section{Implementation}
\label{sec:impl}
The system has three small components, all of which we release.

\emph{Detector + control loop.} A userspace agent ($\approx$40 lines) polls the node-local
ztunnel's saturation signal once per second---worker busy-rate from the Tokio runtime
metrics endpoint, and CPU-throttling counts from the pod's cgroup---and drives a
three-state machine (\textsc{open} $\to$ \textsc{paced} $\to$ \textsc{open}) with a high/low
threshold band and a dwell timer. It emits the target new-connection rate as the only
actuation output. Because ztunnel is distroless, the agent reads the cgroup counters from
the node filesystem and scrapes the metrics port over the pod network, requiring no changes
to ztunnel.

\emph{Data-path enforcer.} The enforcer applies the target rate to the aggressor's
new connections. We implement it two ways behind the same interface (\S\ref{sec:setup-ebpf}):
(i) a connection-admission token bucket, used for all closed-loop experiments; and (ii) a
$\approx$50-line eBPF program attached to the pod's virtual interface via TC that
rate-limits connection-initiating (SYN) packets using a token bucket held in a BPF map the
control loop updates; non-SYN packets take an early return. Both enforce the identical
new-connection-rate policy.

\emph{No mesh or application changes.} The design is external to both Istio and the
application: it reads existing ztunnel telemetry and acts at the CNI/connection layer. It
requires no modification to ztunnel, to the mesh control plane, or to any workload, which
is what makes it deployable on an unmodified ambient-mode cluster.

\section{Why Existing Mechanisms Do Not Address This}
\label{sec:diff}
\paragraph{Byte-pacing is blind to churn.} The closest existing mechanism is Cilium's
Bandwidth Manager, which paces a pod's egress \emph{bytes} using the same kind of eBPF
data-path machinery we use. If byte-pacing could solve this, our contribution would be
redundant. It cannot. Figure~\ref{fig:diff} runs the same churn attack under three
mitigations. At a $1\,$Mbit cap the victim tail is unchanged from unmitigated
($56\,$ms). Even at an aggressive $256\,$kbit cap it only drops to a noisy $29\,$ms with
the proxy still saturated---and such a cap would cripple any legitimate byte-heavy
traffic from that pod. The reason is exactly the mismatch from \S\ref{sec:char}: churn is
cheap on bytes ($\sim$1.25\,MB/s of tiny requests) but expensive on connections, so
thousands of connections slip under any sane byte budget. Our connection-rate controller,
on identical hardware and attack, restores the victim to $3.5\,$ms.

\begin{figure}[t]
\centering
\includegraphics[width=0.78\linewidth]{figs/fig_differentiation.pdf}
\caption{Same attack, three mitigations ($n{=}3$, $\pm$1 SD). Byte-rate pacing (Cilium
Bandwidth Manager) barely helps because it controls the wrong quantity; pacing
new-connection rate restores the victim.}
\label{fig:diff}
\end{figure}

\paragraph{Autoscaling does not apply.} A natural response is to scale out. But
ztunnel is a fixed one-per-node DaemonSet: there is no ``add more ztunnels'' knob, so the
usual horizontal escape hatch simply does not exist for the shared proxy. Scaling the
\emph{victim's} application also does not help, because the victim's app is not the
bottleneck (its CPU is idle; \S\ref{sec:char}). And even where autoscaling could act, its
seconds-scale reaction time cannot catch a saturation onset that, as \S\ref{sec:char}
showed, is effectively instantaneous.

\section{Overhead and Sensitivity}
\label{sec:overhead}
The mitigation is inexpensive and safe to leave enabled.
The data-path cost is dominated by the eBPF enforcer, which we microbenchmark in
isolation at $9\,$ns per packet (mean over $10^6$ iterations) and which runs \emph{only}
on connection-initiating (SYN) packets---steady data traffic hits an early return and is
untouched. We report this as the per-packet cost of the enforcement primitive, not as an
end-to-end overhead measurement; the userspace control loop adds one metric scrape and a
map write per second, which is negligible and off the data path. When the proxy is
healthy the controller holds the cap open, so the \emph{whole system} is a measured
no-op: with the controller running but no aggressor, victim P99 ($2.9\,$ms) is
statistically identical to having no controller at all---this end-to-end no-op result
uses the same connection-layer backend as our headline numbers, so it reflects the
deployed cost, not just the microbenchmark. The setpoint is forgiving rather than a knife-edge (Figure~\ref{fig:sens})---the
victim stays protected ($<4.5\,$ms) across a $6\times$ range of connection-rate caps
(200--1200 conns/s), degrading gracefully as the cap approaches the proxy's knee. And on
controller failure the last cap persists (fail-safe direction) rather than reverting to
unlimited.

\begin{figure}[t]
\centering
\includegraphics[width=0.78\linewidth]{figs/fig_sensitivity.pdf}
\caption{The controller's setpoint is forgiving: the victim stays well-protected across a
$6\times$ range of connection-rate caps (shaded), rather than depending on a precisely
tuned value.}
\label{fig:sens}
\end{figure}

\section{Threat Model and Robustness}
\label{sec:threat}
We separate what our design defends against today from where an informed adversary could
push, since ``noisy neighbor'' spans both accidental and adversarial workloads.

\paragraph{What we defend against.} A co-located tenant---malicious or merely
misconfigured---that saturates the node's shared proxy by opening connections faster than
the proxy can service them. This is the common accidental case (a client library with
aggressive reconnection, short keep-alive, or per-request TLS) and the simplest
deliberate one. Our controller detects the resulting proxy saturation and paces the
offending source's new-connection rate back below the knee.

\paragraph{Evasion 1: staying just under the threshold.} An adversary who knows the
detector could try to hover just below the saturation point. This is self-limiting by
construction: the threshold is the proxy's saturation knee, so ``just under'' is, by
definition, the regime where the victim is \emph{not} harmed (Figure~\ref{fig:principle}
shows the victim is at baseline for any load the proxy can absorb). The adversary cannot
both stay below the knee and hurt the neighbor---the knee is exactly where harm begins.
The residual risk is threshold \emph{tuning}: set too high, a slow ramp could inflict some
tail before tripping; our hysteretic design trips at the first saturation sample
($\sim$1\,s), and an adaptive setpoint (below) tightens this further.

\paragraph{Evasion 2: spreading churn across many source IPs.} A stronger adversary could
churn connections from many pods/IPs so that no single source looks abusive. This is a
real limitation of a purely per-source cap and we do not claim to solve it here. Two
things bound it. First, the \emph{detection} signal is source-agnostic---it is the
proxy's own saturation---so the system still \emph{knows} it is under attack even when no
single source stands out; what degrades is the \emph{attribution} needed to pace fairly.
Second, the natural extension is to pace against the proxy's total connection-admission
budget and divide it fairly across active sources (a weighted fair share, \S\ref{sec:disc}),
which caps aggregate churn regardless of how many identities it is spread over. We
evaluate the single-source case here and scope multi-source fair-sharing as the primary
piece of future work.

\paragraph{Multi-tenant reality.} Production nodes host many pods, not one victim and one
aggressor. Our detection generalizes unchanged (the signal is the shared proxy's state),
but the actuation policy must move from ``cap the aggressor'' to ``allocate the proxy's
connection budget across sources.'' We sketch this in \S\ref{sec:disc} but do not yet
evaluate it; it is the most important open item for a deployment.

\section{Discussion, Limitations, and Future Work}
\label{sec:disc}
\paragraph{The general lesson.} The recurring theme is a resource--control mismatch.
Shared-dataplane designs concentrate a new scarce resource---proxy connection-handling
capacity---in a component that (a) no application owns, (b) cannot be scaled out, and
(c) is invisible to byte- and CPU-based tools. The fix is not clever; it is just aimed at
the right target: measure the scarce resource where it is scarce (the proxy), and control
the input that consumes it (new connections). We expect the same recipe to apply to other
shared-dataplane proxies.

\paragraph{Honest limitations.} We group these by what closes them.
\emph{(Statistical.)} The headline mitigation result is over $n{=}12$ trials with 95\%
bootstrap CIs and a significance test (\S\ref{sec:mitigation}). The differentiation
(\S\ref{sec:diff}) and sensitivity (\S\ref{sec:overhead}) experiments remain smaller
($n{=}3$ and $n{=}1$--$3$ per point); because latency is heavy-tailed, a camera-ready
should bring these to the same $n{=}12$ standard and lengthen the graceful-degradation
soak. The effect sizes there are large enough (byte-pacing does not move the tail at all;
the setpoint is protective across a $6\times$ range) that we expect the conclusions to
hold, but we report them as the less-powered results they are.
\emph{(Integration.)} The eBPF enforcer is microbenchmarked, not yet wired ahead of
Cilium's tcx hook (\S\ref{sec:setup-ebpf}); the closed-loop numbers use the equivalent
connection-admission backend, and finishing the in-CNI attach is engineering, not new
mechanism.
\emph{(Generality.)} We characterized one victim endpoint, one benchmark
(DeathStarBench Social Network), and one mesh (Istio ambient). We argue the mechanism
(connection-setup cost) is transport-agnostic, but this is currently reasoning, not
measurement: the L7 waypoint proxy, a second (e.g., gRPC) workload, and another
shared-dataplane mesh are untested and are the clearest way to substantiate the generality
claim. \emph{(Multi-tenant \& adversarial.)} Covered in \S\ref{sec:threat}: single-source
today; source-agnostic \emph{detection} already generalizes, but fair \emph{actuation}
across many simultaneous churn sources is the primary open problem.

\paragraph{Adaptivity.} Our controller uses a fixed setpoint chosen from the measured
knee. A production version should learn it online---tighten the cap until throttling
stops, then hold just below the knee---and split capacity fairly when several sources
contend. The mechanism classifier (\S\ref{sec:signal}) also suggests routing: apply
connection pacing when throttling indicates a proxy problem, and escalate to
backend/app remediation when the fingerprint indicates the other cause.

\section{Related Work}
\paragraph{Overload control and admission control.} Breakwater and DAGOR protect a server
under load, and Aequitas does admission control for datacenter RPCs. All three act by
\emph{dropping or deferring requests}, keyed off queueing delay or per-RPC latency, at the
RPC or transport layer. They are the right conceptual comparison, and two points relate
them to our setting. First, their \emph{signal} is request-level congestion at an
endpoint the application controls; ours is compute saturation of a \emph{shared network
proxy} that no application owns and that request-level metrics do not observe (in our
backend-contention case the proxy is busy holding tasks; in the churn case the cost is in
connection setup, not request queueing). Second, their \emph{lever} is request
admission---which, adapted to our problem, is essentially per-connection admission at the
proxy, i.e., a proxy-resident fair queue over connections. That is a legitimate
alternative enforcement point to our CNI-layer pacing, and comparing CNI-layer connection
pacing against a proxy-resident connection fair-share (a ztunnel-internal Breakwater-style
admission controller) is the baseline we most want to add next; it needs modifying ztunnel
itself, which is why it is future rather than present work. The key distinction is not
pace-vs-drop alone but \emph{where the scarce resource is and what signal exposes it}:
we show the shared-proxy connection-handling bottleneck is invisible to the byte- and
request-level signals these systems use.

\paragraph{Pacing mechanism and CNI QoS.} EDT-based pacing itself is not new---it is the
mechanism behind Carousel and Cilium's Bandwidth Manager---but those pace \emph{bytes}
from static per-pod limits; our contribution is closing a loop from proxy saturation onto
\emph{connection-rate} pacing, and showing empirically (\S\ref{sec:diff}) that the
byte-rate incumbent cannot address connection-churn interference. Interference-aware
schedulers reduce co-location contention by managing CPU allocation; we act in the network
data path on a resource (proxy connection handling) that scheduling does not govern.

\paragraph{Positioning.} To our knowledge this is the first measurement and mitigation of
connection-churn interference through a shared ambient-mesh proxy. Our novelty is not a new
pacing primitive but the diagnosis (what is scarce, what signal exposes it) and the
matching closed loop.

\section{Conclusion}
Sharing a proxy to save resources creates a shared point of failure whose scarce
resource---connection handling---is invisible to the tools operators already run. We
measured this effect on a real microservice benchmark, separated it from a look-alike
backend-contention cause, showed that the proxy's own telemetry detects and classifies it
almost perfectly (though it neither predicts severity nor warns early), and built a small
reactive controller that restores victim tail latency $12.7\times$ where byte-rate pacing
cannot. The fix is simple once the diagnosis is right: watch the resource that is truly
scarce, and pace the input that consumes it.

\end{document}
